% !TEX program = xelatex
\documentclass[12pt,a4paper]{ctexart}
\usepackage{amsmath,amssymb}
\usepackage{pifont}
\newcommand{\cnum}[1]{\ding{\numexpr171+#1\relax}}
\usepackage[margin=1.8cm]{geometry}
\usepackage{enumitem}
\usepackage{tabularx}
\usepackage{xcolor}
\usepackage{fancyhdr}

\pagestyle{fancy}
\fancyhf{}
\renewcommand{\headrulewidth}{0pt}
\renewcommand{\footrulewidth}{0pt}
\fancyfoot[R]{\thepage}
\fancypagestyle{plain}{%
  \fancyhf{}%
  \renewcommand{\headrulewidth}{0pt}%
  \fancyfoot[R]{\thepage}%
}

\setlist[enumerate]{itemsep=1.1em, parsep=0pt, topsep=0.5em, leftmargin=2.4em}
\setlength{\parskip}{0pt}
\setlength{\parindent}{0pt}
\linespread{1.15}
\renewcommand{\arraystretch}{1.2}
\setlength{\extrarowheight}{2pt}
\setlength{\lineskiplimit}{3pt}
\setlength{\lineskip}{4pt plus 2pt}

% 选项排版：ABCD 固定列位，按内容长度选 4 列 / 2 列 / 1 列
\newcommand{\optfour}[4]{\par\begin{tabularx}{\linewidth}{@{}X@{\quad}X@{\quad}X@{\quad}X@{}}#1&#2&#3&#4\end{tabularx}}
\newcommand{\opttwo}[4]{\par\begin{tabularx}{\linewidth}{@{}X@{\quad}X@{}}#1&#2\\[0.8em]#3&#4\end{tabularx}}
\newcommand{\optone}[4]{\par\begin{tabularx}{\linewidth}{@{}X@{}}#1\\[0.8em]#2\\[0.8em]#3\\[0.8em]#4\end{tabularx}}

\begin{document}

\begin{center}
{\Large\bfseries 概统第三讲《多维随机变量及其分布》基础诊断卷（30 题）}
\end{center}

\subsection*{一、客观题（1--20 题）}

\begin{enumerate}[label=\arabic*.]

\item（帕概-C-1-例3.1 / 姜基础）下列四个函数中，\textbf{不能}作为二维随机变量联合分布函数的是（\quad）
\optone{(A) $F_{1}(x,y) = \begin{cases} (1 - \mathrm{e}^{-x})(1 - \mathrm{e}^{-y}), & x > 0, y > 0, \\ 0, & \text{其他} \end{cases}$}%
{(B) $F_{2}(x,y) = \dfrac{1}{\pi^{2}}\left(\dfrac{\pi}{2} + \arctan\dfrac{x}{2}\right)\left(\dfrac{\pi}{2} + \arctan\dfrac{y}{3}\right)$}%
{(C) $F_{3}(x,y) = \begin{cases} 1, & x + 2y \geqslant 1, \\ 0, & x + 2y < 1 \end{cases}$}%
{(D) $F_{4}(x,y) = \begin{cases} 1 - 2^{-x} - 2^{-y} + 2^{-x-y}, & x > 0, y > 0, \\ 0, & \text{其他} \end{cases}$}

\item（李林880-基-18-选1）设二维随机变量 $(X, Y)$ 的分布函数为 $F(x, y)$，则概率 $P\{X > x_{0}, Y > y_{0}\} = (\quad)$
\opttwo{(A) $1 + F(x_{0}, y_{0}) - F(x_{0}, +\infty) - F(+\infty, y_{0})$}{(B) $F(x_{0}, y_{0}) - 1 + F(x_{0}, +\infty) + F(+\infty, y_{0})$}{(C) $1 - F(x_{0}, +\infty) - F(+\infty, y_{0})$}{(D) $1 - F(x_{0}, y_{0})$}

\item（帕概-C-2.a-1990/1997数三）设随机变量 $X$ 和 $Y$ 相互独立且同分布，其概率分布均为：
\[
P\{X = -1\} = P\{Y = -1\} = \frac{1}{2}, \quad P\{X = 1\} = P\{Y = 1\} = \frac{1}{2},
\]
则下列各式中正确的是（\quad）
\optfour{(A) $X = Y$}{(B) $P\{X = Y\} = 0$}{(C) $P\{X = Y\} = \dfrac{1}{2}$}{(D) $P\{X + Y = 0\} = \dfrac{1}{4}$}

\item（帕概-C-2.a-2005数一三）设二维随机变量 $(X, Y)$ 的联合分布律为：
\begin{center}
\begin{tabular}{|c|c|c|}
\hline
$X \backslash Y$ & $0$ & $1$ \\ \hline
$0$ & $0.4$ & $a$ \\ \hline
$1$ & $b$ & $0.1$ \\ \hline
\end{tabular}
\end{center}
已知随机事件 $\{X = 0\}$ 与 $\{X + Y = 1\}$ 相互独立，则（\quad）
\optfour{(A) $a = 0.2, b = 0.3$}{(B) $a = 0.4, b = 0.1$}{(C) $a = 0.3, b = 0.2$}{(D) $a = 0.1, b = 0.4$}

\item（张千-强-概4-T3 / 数一交叉）设二维随机变量 $(X, Y)$ 的概率密度为：
\[
f(x, y) = \begin{cases} \dfrac{1}{4}, & -1 \leqslant x < 1, 0 \leqslant y < 2, \\ 0, & \text{其他}, \end{cases}
\]
则二次型 $g(x_1, x_2, x_3) = x_1^2 + 2x_2^2 + Yx_3^2 + 2x_1x_2 + 2Xx_1x_3$ 正定的概率为（\quad）
\optfour{(A) $\dfrac{2}{3}$}{(B) $\dfrac{1}{2}$}{(C) $\dfrac{1}{3}$}{(D) $\dfrac{1}{4}$}

\item（帕概-C-2.a-2023数一）设随机变量 $X$ 与 $Y$ 相互独立，且 $X \sim B\left(1, \dfrac{1}{3}\right)$，$Y \sim B\left(2, \dfrac{1}{2}\right)$，则概率 $P\{X = Y\} = (\quad)$
\optfour{(A) $\dfrac{1}{6}$}{(B) $\dfrac{1}{3}$}{(C) $\dfrac{1}{2}$}{(D) $\dfrac{5}{12}$}

\item（帕概-C-3.a(2)-2012数一）设随机变量 $X$ 与 $Y$ 相互独立，且分别服从参数为 $1$ 与参数为 $4$ 的指数分布，则概率 $P\{X < Y\} = (\quad)$
\optfour{(A) $\dfrac{1}{5}$}{(B) $\dfrac{1}{3}$}{(C) $\dfrac{2}{3}$}{(D) $\dfrac{4}{5}$}

\item（帕概-C-3.b(1)-1998数一）设平面区域 $D$ 由曲线 $y = \dfrac{1}{x}$ 及直线 $y = 0, x = 1, x = \mathrm{e}^{2}$ 所围成。二维随机变量 $(X, Y)$ 在区域 $D$ 上服从均匀分布，则 $(X, Y)$ 关于 $X$ 的边缘概率密度在 $x = 2$ 处的值 $f_{X}(2) = (\quad)$
\optfour{(A) $\dfrac{1}{2}$}{(B) $\dfrac{1}{4}$}{(C) $\dfrac{1}{8}$}{(D) $\dfrac{1}{2\mathrm{e}^{2}}$}

\item（帕概-C-3.b(1)-1992数三）设二维随机变量 $(X, Y)$ 的概率密度为：
\[
f(x, y) = \begin{cases} \mathrm{e}^{-y}, & 0 < x < y, \\ 0, & \text{其他}. \end{cases}
\]
则下列断言正确的是（\quad）
\opttwo{(A) $X$ 与 $Y$ 相互独立，且 $X \sim E(1)$}{(B) $X$ 与 $Y$ 相互独立，且 $Y \sim E(1)$}{(C) $X$ 与 $Y$ 不相互独立，且 $X$ 的边缘密度在 $x > 0$ 时为 $f_{X}(x) = \mathrm{e}^{-x}$}{(D) $X$ 与 $Y$ 不相互独立，且 $X$ 的边缘密度在 $x > 0$ 时为 $f_{X}(x) = x\mathrm{e}^{-x}$}

\item（帕概-C-4-2007数一三）设随机变量 $(X, Y)$ 服从二维正态分布，且 $X$ 与 $Y$ 不相关，$f_{X}(x), f_{Y}(y)$ 分别表示 $X, Y$ 的边缘概率密度，则在 $Y = y$ 的条件下，$X$ 的条件概率密度 $f_{X|Y}(x \mid y)$ 为（\quad）
\optfour{(A) $f_{X}(x)$}{(B) $f_{Y}(y)$}{(C) $f_{X}(x)f_{Y}(y)$}{(D) $\dfrac{f_{X}(x)}{f_{Y}(y)}$}

\item（帕概-C-4-1999数一）设两个相互独立的随机变量 $X$ 和 $Y$ 分别服从正态分布 $N(0, 1)$ 和 $N(1, 1)$，则必有（\quad）
\opttwo{(A) $P\{X + Y \leqslant 0\} = \dfrac{1}{2}$}{(B) $P\{X + Y \leqslant 1\} = \dfrac{1}{2}$}{(C) $P\{X - Y \leqslant 0\} = \dfrac{1}{2}$}{(D) $P\{X - Y \leqslant 1\} = \dfrac{1}{2}$}

\item（帕概-C-4-2025数一）设二维随机变量 $(X, Y)$ 服从二维正态分布 $N(0, 0; 1, 1; \rho)$，其中相关系数 $\rho \in (-1, 1)$。若实数 $a, b$ 满足 $a^{2} + b^{2} = 1$，则方差 $D(aX + bY)$ 的最大值为（\quad）
\optfour{(A) $1$}{(B) $2$}{(C) $1 + |\rho|$}{(D) $1 + \rho^{2}$}

\item（帕概-C-4-2020数三）设随机变量 $(X, Y)$ 服从二维正态分布 $N\left(0, 0; 1, 4; -\dfrac{1}{2}\right)$，则下列随机变量中服从标准正态分布且与 $X$ 相互独立的是（\quad）
\opttwo{(A) $\dfrac{\sqrt{5}}{5}(X + Y)$}{(B) $\dfrac{\sqrt{5}}{5}(X - Y)$}{(C) $\dfrac{\sqrt{3}}{3}(X + Y)$}{(D) $\dfrac{\sqrt{3}}{3}(X - Y)$}

\item（帕概-C-4-2022数一）设随机变量 $X \sim N(0, 1)$。若在给定 $X = x$ 的条件下，随机变量 $Y \sim N(x, 1)$，则 $X$ 与 $Y$ 的相关系数 $\rho_{XY}$ 为（\quad）
\optfour{(A) $\dfrac{1}{4}$}{(B) $\dfrac{1}{2}$}{(C) $\dfrac{\sqrt{3}}{3}$}{(D) $\dfrac{\sqrt{2}}{2}$}

\item（张千-强-概5-T1）设 $X_{1}, X_{2}$ 是来自标准正态总体 $N(0, 1)$ 的简单随机样本，则商变量 $Y = \dfrac{X_{1}}{X_{2}}$ 的概率密度 $f_{Y}(y) = (\quad)$
\optfour{(A) $\dfrac{1}{\pi(1 + y^{2})}$}{(B) $\dfrac{1}{\pi(1 + y)}$}{(C) $\dfrac{1}{1 + y^{2}}$}{(D) $\dfrac{1}{\sqrt{\pi}} \mathrm{e}^{-y^{2}}$}

\item（帕概-C-5.c-2009数一三）设随机变量 $X$ 与 $Y$ 相互独立，且 $X$ 服从标准正态分布 $N(0, 1)$，$Y$ 的概率分布为 $P\{Y = 0\} = P\{Y = 1\} = \dfrac{1}{2}$。记 $F_{Z}(z)$ 为随机变量 $Z = XY$ 的分布函数，则函数 $F_{Z}(z)$ 的间断点个数为（\quad）
\optfour{(A) $0$}{(B) $1$}{(C) $2$}{(D) $3$}

\item（帕概-C-5.b-2024数一三）设随机变量 $X, Y$ 相互独立，且均服从参数为 $\lambda\ (\lambda > 0)$ 的指数分布。令 $Z = |X - Y|$，则下列随机变量中与 $Z$ 同分布的是（\quad）
\optfour{(A) $X + Y$}{(B) $\dfrac{X + Y}{2}$}{(C) $2X$}{(D) $X$}

\item（张千-强-概5-T9）设二维随机变量 $(X, Y)$ 在区域 $D = \{(x, y) \mid |x| + |y| \leqslant 1\}$ 上服从均匀分布，令 $U = X + Y, V = X - Y$，则协方差 $\mathrm{Cov}(U, V) = $ \underline{\hspace{5em}}。

\item（李林880-基-18-填3）设随机变量 $X, Y$ 均服从区间 $[0, 4]$ 上的均匀分布且相互独立。若 $P\{\max\{X, Y\} \leqslant 3\} = \dfrac{9}{16}$，则概率 $P\{\min\{X, Y\} > 3\} = $ \underline{\hspace{5em}}。

\item（帕概-C-4-2011数一三）设二维随机变量 $(X, Y)$ 服从二维正态分布 $N(\mu, \mu; \sigma^{2}, \sigma^{2}; 0)$，则数学期望 $E(XY^{2}) = $ \underline{\hspace{5em}}。

\end{enumerate}

\subsection*{二、解答题（21--30 题）}

\setlist[enumerate]{itemsep=0.6em, parsep=0pt, topsep=0.3em, leftmargin=2.4em}

\begin{enumerate}[label=\arabic*.,resume]

\item（帕概-C-1-1995数三）已知随机变量 $X$ 和 $Y$ 的联合概率密度为：
\[
\varphi(x, y) = \begin{cases} 4xy, & 0 \leqslant x \leqslant 1, 0 \leqslant y \leqslant 1, \\ 0, & \text{其他}. \end{cases}
\]
求 $X$ 和 $Y$ 的联合分布函数 $F(x, y)$。

\item（帕概-C-2.b-2020数三）设二维随机变量 $(X, Y)$ 在区域 $D = \left\{(x, y) \;\middle|\; 0 < y < \sqrt{1 - x^{2}}\right\}$ 上服从均匀分布。令指示变量：
\[
Z_{1} = \begin{cases} 1, & X - Y > 0, \\ 0, & X - Y \leqslant 0, \end{cases} \quad Z_{2} = \begin{cases} 1, & X + Y > 0, \\ 0, & X + Y \leqslant 0. \end{cases}
\]
(1) 求二维离散型随机变量 $(Z_{1}, Z_{2})$ 的联合概率分布；\\
(2) 求 $Z_{1}$ 与 $Z_{2}$ 的相关系数 $\rho$。

\item（帕概-C-2.c-2001数一）设某班车起点站上车人数 $X$ 服从参数为 $\lambda\ (\lambda > 0)$ 的泊松分布，每位乘客在中途下车的概率为 $p\ (0 < p < 1)$，且各乘客中途下车与否相互独立。以 $Y$ 表示在中途下车的乘客人数。\\
(1) 在发车时有 $n$ 个乘客（即 $X = n$）的条件下，求中途有 $m$ 人下车的条件概率 $P\{Y = m \mid X = n\}\ (m = 0, 1, \dots, n)$；\\
(2) 求二维随机变量 $(X, Y)$ 的联合分布律 $P\{X = n, Y = m\}$；\\
(3) 求随机变量 $Y$ 的边缘分布律，并指出 $Y$ 服从什么概率分布及其参数。

\item（帕概-C-3.b(2)-例3.4 / 姜强化）设随机变量 $X \sim U(0, 1)$，在给定 $X = x\ (0 < x < 1)$ 的条件下，随机变量 $Y \sim U(x, 1)$。\\
(1) 求 $(X, Y)$ 的联合概率密度 $f(x, y)$；\\
(2) 求 $(X, Y)$ 关于 $Y$ 的边缘概率密度 $f_{Y}(y)$；\\
(3) 计算概率 $P\{X + Y > 1\}$。

\item（李林880-基-18-大9 / 正态不相关非独立经典反例）设二维随机变量 $(X, Y)$ 的概率密度为：
\[
f(x, y) = \frac{1}{k} (1 + xy)\mathrm{e}^{-\frac{1}{2}(x^{2} + y^{2})}, \quad k > 0, (x, y) \in \mathbb{R}^{2}.
\]
(1) 求常数 $k$ 的值；\\
(2) 判别 $X$ 与 $Y$ 是否不相关，并严格判定 $X$ 与 $Y$ 是否相互独立；\\
(3) 若随机变量 $Z$ 服从区间 $[-\pi, \pi]$ 上的均匀分布，且 $X$ 与 $Z$ 相互独立，求 $U = X + Z$ 的概率密度（可用标准正态分布函数 $\Phi(x)$ 表示）。

\newpage
\item（帕概-C-4-1994数一）已知二维随机变量 $(X, Y)$ 服从二维正态分布，且 $X \sim N(1, 3^{2})$，$Y \sim N(0, 4^{2})$，$X$ 与 $Y$ 的相关系数 $\rho_{XY} = -\dfrac{1}{2}$。令随机变量：
\[
Z = \frac{X}{3} + \frac{Y}{2}.
\]
(1) 求 $Z$ 的数学期望 $E(Z)$ 和方差 $D(Z)$；\\
(2) 求 $X$ 与 $Z$ 的协方差 $\mathrm{Cov}(X, Z)$ 及相关系数 $\rho_{XZ}$；\\
(3) 问随机变量 $X$ 与 $Z$ 是否相互独立？写出严密的判定依据与证明。

\item（帕概-C-5.b-2007数一二三）设二维随机变量 $(X, Y)$ 的概率密度为：
\[
f(x, y) = \begin{cases} 2 - x - y, & 0 < x < 1, 0 < y < 1, \\ 0, & \text{其他}. \end{cases}
\]
(1) 计算概率 $P\{X > 2Y\}$；\\
(2) 求随机变量 $Z = X + Y$ 的概率密度 $f_{Z}(z)$。

\item（帕概-C-5.b-2023数一）设二维随机变量 $(X, Y)$ 的概率密度为：
\[
f(x, y) = \begin{cases} \dfrac{2}{\pi}(x^{2} + y^{2}), & x^{2} + y^{2} \leqslant 1, \\ 0, & \text{其他}. \end{cases}
\]
(1) 求 $X$ 与 $Y$ 的协方差 $\mathrm{Cov}(X, Y)$；\\
(2) 问 $X$ 与 $Y$ 是否相互独立？说明理由；\\
(3) 求随机变量 $Z = X^{2} + Y^{2}$ 的概率密度 $f_{Z}(z)$。

\item（李林880-综-18-大8 / 微分方程极坐标交叉）设随机变量 $X$ 与 $Y$ 相互独立且服从同一分布，当 $x \leqslant 0$ 时，$X$ 的概率密度 $f(x) = 0$；当 $x > 0$ 时，$f(x)$ 满足微分方程：
\[
f'(x) + 2x f(x) = 0.
\]
令随机变量 $Z = \sqrt{X^{2} + Y^{2}}$。\\
(1) 求 $X$ 的概率密度 $f(x)$；\\
(2) 求随机变量 $Z$ 的分布函数 $F_{Z}(z)$ 与概率密度 $f_{Z}(z)$；\\
(3) 计算方差 $D(Z)$。

\item（李林880-综-18-大7 / 概率积分变换+连续卷积压轴）设随机变量 $X$ 与 $Y$ 相互独立，其中 $X$ 的概率密度为：
\[
f(x) = \begin{cases} x, & 0 \leqslant x < 1, \\ 2 - x, & 1 \leqslant x < 2, \\ 0, & \text{其他}, \end{cases}
\]
$Y$ 服从参数为 $2$ 的指数分布 $Y \sim E(2)$。$F(x)$ 为 $X$ 的分布函数，记随机变量 $Z = F(X) + Y$。\\
(1) 求随机变量 $F(X)$ 的分布函数与概率密度；\\
(2) 求随机变量 $Z$ 的概率密度 $f_{Z}(z)$。

\end{enumerate}

\end{document}
